\section{Inputs}\label{app:features}

\begin{table}[H]\centering\small
\caption{Model inputs, measured at December of the snapshot year.}\label{tab:features}
\begin{tabularx}{\linewidth}{>{\raggedright\arraybackslash}p{5.2cm}X}\toprule
Input & Definition \\ \midrule
Credit score & Borrower credit score at origination (missing if not disclosed) \\
Original LTV, CLTV & Loan-to-value and combined loan-to-value at origination (\%) \\
Debt-to-income & Back-end ratio at origination (\%; missing above 65) \\
Original rate; spread & Note rate; note rate minus the PMMS 30-year rate in the origination month \\
Purpose, occupancy, property type, channel, first-time buyer & Categorical, as disclosed \\
Units; number of borrowers & As disclosed \\
MI percentage & Mortgage-insurance coverage at origination \\
Original term & Months \\
Loan age & Months since the first payment date, plus one \\
Current LTV & Current balance over (original balance / original LTV) $\times$ state HPI now / state HPI at origination \\
HPI growth & State FHFA index, fourth quarter over fourth quarter \\
Unemployment; change & State rate in December; change over twelve months \\
\bottomrule
\end{tabularx}
\end{table}

\section{Data log, deviations, corrections and added analyses}\label{app:deviations}

The design was committed at 10:36 ET on 29 October 2026 (commit
\texttt{07c8ffd}); the year archives were downloaded between 13:25 and 14:30
ET the same day, and the first loan-level file was opened at 14:33 ET to read
the Release 47 layout. The commit timestamp is the author's machine time; it
is not a third-party registration.

\paragraph{Deviation 1: prepayment grid.} The design chose the prepayment
model as the best of the full 65-model grid. A timing probe at December 2008,
before any allowance had been computed, showed that a prepayment LightGBM
model runs to about 2,700 boosting rounds (about 100 seconds), so the full grid
would have cost about two hours per date. The prepayment model is instead the
best of the eight unconstrained LightGBM settings at seed 1 and the spline
logit. The design's monotone constraints are written for default risk and
would force prepayment to fall as the credit score rises; extra seeds only
choose the luckiest draw of a model the design holds fixed. The grid of
default models, which is what the study is about, is unchanged.

\paragraph{Implementation details the design did not specify.} Early stopping
uses a further 10\% of loans (so fitting uses 70\%). A final part-year hazard
is scaled as $1-(1-h)^{q/12}$. The LGD table falls back to the band, then the
overall, mean for cells with fewer than 30 resolutions. A default later
repaid voluntarily is a zero-loss resolution; defaults removed by a
reperforming sale or a defect, or with no disclosed loss, are not used in the
severity table. The IFRS 9 origination projection is defined in
\cref{sec:predictions}. All of these were fixed before any allowance was
computed.

\paragraph{Deviations 2 to 4.} (2) The design listed a secondary analysis in
which severity varies with the default model; it was not run, and severity is
held fixed throughout. (3) The design scales the scored sample's allowance by
the inverse sampling fraction; the code uses the ratio estimator (expected
loss over balance on the sample, times the balance of the whole active
book), which agrees in expectation and has lower variance. (4) The design
computes the reversion means ``on data before $t$''; the code uses every year
up to and including $T$, whose December values are observed at $T$.

\paragraph{Corrections.} Four bugs were found while checking results and
fixed; the affected steps were rerun for every date, and every number in the
paper comes from the corrected code. The fitted models were not affected by
any of them. (1) The realized macro path looked up each state with the wrong
key and fell back to the long-run means. (2) A few dozen loans a year with no
remaining term were cast to invalid integers; they now use the original term
less the months paid. Both were found after the first three dates'
allowances. (3) The severity table first averaged per-loan ratios of loss to
balance at default, which a handful of loans with near-zero balances at
default dominate (one with \$0.19 at default and a \$27{,}168 loss); severity
is now balance-weighted, total loss over total balance at default. This was
found after every date's allowance existed, from an implausible level.
(4) Found by an independent review: ``immediate'' reversion set the macro
value at year-end $T+2$ to the long-run mean, which, because that value is
the house-price growth over year $T+2$, shortened the two-year forecast
period. Immediate reversion now starts at $T+3$ (\cref{sec:method}).

\paragraph{Analyses added after seeing results.} Each is labelled in the text
and none replaces a recorded test: (1) the paired-difference tolerance
(\cref{sec:sensitivity}); (2) the like-for-like versions of H1 and H6 and
(3) the refits on independent samples, with the accuracy and bias of both
kinds of refit (\cref{sec:refits}); (4) the lower bound for realized losses,
which counts defaults unresolved in June 2026 at zero; (5) the decomposition
of the H5 miss (\cref{tab:decomp}); (6) the delinquency of the portfolio
(\cref{sec:limits}); (7) the position of the most accurate model in the set,
the rank correlation of accuracy with the allowance, H4 on equal-size sets,
the refit ranges against realized losses, and the decomposition with the
December 2019 severity table and the realized default rates
(\texttt{code/07b\_posthoc.py}). Items (3), (5) and (6), and bug (4),
followed an independent referee review of a complete draft, and item (7) a
second one.

\section{Severity tables}\label{app:lgd}

\begin{table}[H]\centering\footnotesize
\caption{Loss given default (\% of balance at default) by current-LTV band at the start of the default year; P = power of sale, J = judicial. $^\dagger$ fewer than 30 resolutions in the cell, so the band or overall mean is used.}\label{tab:lgd}
\setlength{\tabcolsep}{3.5pt}
\begin{tabular}{rrrrrrrrrrrrr}\toprule
Dec. & \multicolumn{2}{c}{$<$60} & \multicolumn{2}{c}{60--80} & \multicolumn{2}{c}{80--90} & \multicolumn{2}{c}{90--100} & \multicolumn{2}{c}{100--120} & \multicolumn{2}{c}{$\ge$120} \\
 & P & J & P & J & P & J & P & J & P & J & P & J \\ \midrule
2006 & 3 & 6 & 11 & 17 & 7 & 12 & 5 & 8 & 10$^\dagger$ & 10$^\dagger$ & 10$^\dagger$ & 10$^\dagger$ \\
2008 & 7 & 9 & 17 & 22 & 12 & 15 & 9 & 10 & 20 & 20 & 15$^\dagger$ & 15$^\dagger$ \\
2012 & 15 & 20 & 31 & 38 & 39 & 42 & 40 & 41 & 48 & 52 & 46 & 47 \\
2017 & 14 & 24 & 28 & 41 & 35 & 47 & 36 & 45 & 41 & 52 & 42 & 48 \\
2025 & 7 & 14 & 18 & 28 & 26 & 37 & 30 & 39 & 39 & 49 & 40 & 47 \\
\bottomrule\end{tabular}
\end{table}


\section{Results by reporting date}\label{app:bydate}

\begin{table}[H]\centering\footnotesize
\caption{Allowance by reporting date. Rates are lifetime expected loss as a percentage of balance on the scored sample; set range in \% of balance; refit and policy widths in basis points of balance; \$ range scales the set's range to the whole active book. Realized loss in parentheses where fewer than eight years are observed.}\label{tab:bydate}
\setlength{\tabcolsep}{4pt}
\begin{tabular}{rrrrrrrrrr}\toprule
& Set & Best & Best & Set range & Rel. & Refit & Policy & \$ range & Realized \\
Dec. & size & AUC & rate (\%) & (\%) & range (\%) & 95\% (bp) & range (bp) & (bn) & (\%) \\ \midrule
2005 & 36 & 0.910 & 0.06 & 0.06--0.09 & 51 & 2.3 & 0.1 & 0.3 & 1.35 \\
2006 & 38 & 0.906 & 0.15 & 0.12--0.18 & 37 & 5.7 & 0.5 & 0.6 & 2.11 \\
2007 & 35 & 0.898 & 0.28 & 0.25--0.38 & 46 & 6.1 & 5.9 & 1.5 & 2.90 \\
2008 & 34 & 0.897 & 0.83 & 0.80--1.76 & 116 & 60.4 & 16.8 & 12.2 & 2.98 \\
2009 & 27 & 0.907 & 1.81 & 1.68--2.69 & 55 & 57.7 & 46.1 & 13.9 & 2.02 \\
2010 & 51 & 0.906 & 1.71 & 1.42--2.68 & 73 & 109.4 & 33.7 & 17.4 & 1.33 \\
2011 & 53 & 0.901 & 2.39 & 2.24--3.57 & 56 & 48.8 & 50.5 & 18.0 & 0.88 \\
2012 & 50 & 0.897 & 1.41 & 1.16--2.37 & 86 & 56.3 & 14.3 & 15.8 & 0.61 \\
2013 & 61 & 0.893 & 0.75 & 0.66--1.01 & 46 & 40.5 & 2.4 & 4.7 & 0.43 \\
2014 & 54 & 0.892 & 0.63 & 0.54--1.01 & 74 & 54.8 & 0.1 & 6.6 & 0.33 \\
2015 & 55 & 0.891 & 0.51 & 0.43--1.11 & 133 & 49.0 & 0.3 & 10.0 & 0.27 \\
2016 & 53 & 0.890 & 0.34 & 0.32--0.63 & 94 & 14.2 & 0.7 & 4.9 & 0.26 \\
2017 & 53 & 0.890 & 0.35 & 0.30--0.48 & 52 & 40.1 & 0.2 & 3.0 & 0.28 \\
2018 & 51 & 0.889 & 0.35 & 0.33--0.67 & 97 & 19.3 & 2.0 & 6.0 & (0.30) \\
2019 & 53 & 0.888 & 0.31 & 0.29--0.40 & 35 & 17.5 & 0.1 & 2.0 & (0.35) \\
2020 & 50 & 0.878 & 0.46 & 0.34--0.79 & 98 & 52.6 & 0.7 & 9.7 & (0.16) \\
2021 & 52 & 0.872 & 0.20 & 0.18--0.25 & 38 & 9.5 & 2.7 & 2.0 & (0.11) \\
2022 & 50 & 0.871 & 0.24 & 0.19--0.42 & 93 & 9.1 & 0.4 & 6.4 & (0.12) \\
2023 & 49 & 0.869 & 0.33 & 0.27--0.36 & 27 & 7.8 & 2.3 & 2.6 & (0.11) \\
2024 & 47 & 0.869 & 0.31 & 0.28--0.39 & 33 & 5.3 & 0.9 & 3.0 & (0.08) \\
2025 & 46 & 0.869 & 0.30 & 0.29--0.36 & 22 & 7.3 & 2.0 & 2.0 & (0.03) \\
\bottomrule\end{tabular}
\end{table}


\section{Data statement}

The study uses only public data and published methods.
